% =====================================================================
%  COMANDOS Y ATAJOS PROPIOS DEL INFORME
%  Objetivo: escribir siempre igual los nombres largos y las unidades,
%  para que el documento quede uniforme aunque lo redacten 5 personas.
% =====================================================================

% --- Tipos de columna para tablas ------------------------------------
% L{3cm} = columna de ancho fijo alineada a la izquierda con salto de línea
% C{3cm} = ídem centrada       R{3cm} = ídem alineada a la derecha
% Se usan sobre todo en la matriz de no conformidades (Tabla 3).
\newcolumntype{L}[1]{>{\raggedright\arraybackslash}p{#1}}
\newcolumntype{C}[1]{>{\centering\arraybackslash}p{#1}}
\newcolumntype{R}[1]{>{\raggedleft\arraybackslash}p{#1}}
% Y = como X de tabularx pero con el texto sin justificar (evita huecos feos)
\newcolumntype{Y}{>{\raggedright\arraybackslash}X}

% --- Nombres recurrentes ---------------------------------------------
% Escribe \mina en lugar de teclear el nombre completo cada vez.
\newcommand{\mina}{Mina Marlin}
\newcommand{\titular}{Montana Exploradora de Guatemala, S.\,A.}
\newcommand{\auditora}{On Common Ground Consultants Inc.}
\newcommand{\hria}{\textit{Human Rights Assessment of Goldcorp's Marlin Mine}}

% --- Unidades y notación técnica -------------------------------------
% Evita que se escriba «g/t» de tres formas distintas en el mismo informe.
\newcommand{\gt}{g/t}                        % ley: gramos por tonelada
\newcommand{\ozAu}{oz Au}                    % onzas de oro
\newcommand{\ozAg}{oz Ag}                    % onzas de plata
\newcommand{\msnm}{m\,s.\,n.\,m.}            % metros sobre el nivel del mar
\newcommand{\tdia}{t/día}                    % toneladas por día

% --- Clasificación de hallazgos (ISO 19011) --------------------------
% Úsalos en la columna «Clasificación» de la matriz de no conformidades
% para que las tres categorías se vean siempre igual.
\newcommand{\NCmayor}{\textbf{NC mayor}}
\newcommand{\NCmenor}{NC menor}
\newcommand{\OBS}{Observación}

% --- Marcador de pendiente -------------------------------------------
% \pendiente{texto} imprime una nota visible en el PDF mientras se redacta.
% CUANDO EL INFORME ESTÉ TERMINADO: comenta la primera definición y
% descomenta la segunda para que los pendientes desaparezcan del PDF final.
\newcommand{\pendiente}[1]{\textcolor{red}{\small[PENDIENTE: #1]}}
%\renewcommand{\pendiente}[1]{}

% --- Cita textual destacada ------------------------------------------
% APA: las citas de más de 40 palabras van en bloque, sin comillas y
% con sangría. \citalarga{texto}{referencia} lo hace automáticamente.
\newenvironment{citalarga}
  {\begin{quote}\singlespacing\small}
  {\end{quote}\normalsize}
